\documentclass[11pt]{article}

\usepackage[utf8]{inputenc}
\usepackage[T1]{fontenc}
\usepackage[margin=1in]{geometry}
\usepackage{amsmath,amssymb}
\usepackage{graphicx}
\usepackage{hyperref}
\usepackage{booktabs}
\usepackage{caption}
\usepackage{subcaption}
\usepackage{array}
\usepackage{multirow}
\usepackage{natbib}
\usepackage{setspace}
\usepackage{xcolor}
\usepackage{url}
\usepackage{textcomp}
\usepackage{float}

\hypersetup{colorlinks=true,linkcolor=blue,urlcolor=blue,citecolor=blue}
\doublespacing

\title{\textbf{A Biologically Informed Explainable Machine Learning Framework for Breast Cancer Progression Risk Stratification From Tumor Gene Expression}}

\author{
Suhaan Thayyil\textsuperscript{1,*} \and Eshaan Nidee\textsuperscript{2} \\[6pt]
\textsuperscript{1}Waxhaw, NC, USA \\
\textsuperscript{2}Frisco, TX, USA \\[4pt]
\textsuperscript{*}Corresponding author: corresponding@email.com
}

\date{February 2026}

\begin{document}

\maketitle

\begin{abstract}
\noindent
Gene expression profiling can supplement traditional breast cancer staging, yet most transcriptomic models sacrifice transparency for performance. We present a framework that collapses RNA-seq expression into seven biologically curated pathway activity scores and evaluates them through both survival analysis and binary classification, with systematic comparisons against standard pathway-scoring baselines and controlled feature ablation.

Seven pathway activity scores---proliferation, estrogen response, immune infiltration, apoptosis, epithelial--mesenchymal transition (EMT), HER2 signaling, and angiogenesis---were computed from RNA-sequencing data in the TCGA-BRCA cohort ($n = 213$) and validated on the SCAN-B GSE96058 cohort ($n = 1{,}483$). Cox proportional hazards models and three binary classifiers (elastic net logistic regression, random forest, gradient boosting) were compared under stratified five-fold cross-validation. SHAP (SHapley Additive exPlanations) analysis provided per-patient explanations. Feature importance stability was assessed via cross-fold rank correlation, and the mean z-score pathway aggregation was benchmarked against a rank-based single-sample enrichment (ssGSEA) baseline.

In survival analysis, the combined pathway-plus-clinical Cox model achieved a cross-validated concordance index (C-index) of 0.827, compared to 0.822 for clinical features alone and 0.626 for pathway features alone. Kaplan--Meier analysis confirmed significant risk-group separation (log-rank $p < 10^{-52}$). In binary classification, the combined model achieved AUC = 0.856 (random forest), compared to 0.849 clinical-only and 0.645 pathway-only. The mean z-score aggregation consistently outperformed the ssGSEA baseline by 0.010--0.038 AUC. SHAP feature rankings were highly stable across cross-validation folds (Spearman $\rho = 0.82$), and the learned attributions recovered well-characterized oncological mechanisms.

These results demonstrate that pathway-oriented explainable machine learning can deliver transparent, generalizable breast cancer risk stratification grounded in recognized cancer biology. \\[6pt]
\noindent\textbf{Keywords:} breast cancer prognosis, gene expression, biological pathway scores, explainable AI, SHAP, Cox survival analysis, feature ablation, TCGA, SCAN-B \\[4pt]
\noindent\textbf{Availability:} Code and data at \url{https://github.com/suhaanthayyil/breast-cancer-prediction}
\end{abstract}


%======================================================================
\section{Introduction}
\label{sec:intro}
%======================================================================

Breast cancer encompasses a family of diseases rather than a single entity. At the molecular level, tumors are commonly sorted into at least four intrinsic subtypes: Luminal A, Luminal B, HER2-enriched, and Basal-like \citep{tcga2012}. Each subtype carries its own expected trajectory, yet within any one category the spread of actual patient outcomes remains wide. Two patients with the same molecular subtype, matched on stage and receptor status, may follow entirely different courses. Subtype labels, despite their proven utility, leave substantial prognostic uncertainty on the table.

Standard clinical variables---tumor size, nodal involvement, histological grade, patient age---explain part of this residual variation but not all of it. RNA sequencing offers a much more granular window into tumor biology, measuring expression across upwards of 20,000 genes simultaneously. This richness comes with a cost: models that ingest thousands of raw gene-level features are vulnerable to overfitting, especially when training cohorts are modest in size, and they produce outputs that are nearly impossible for a treating physician to reason about at the bedside.

\subsection{The Interpretability Gap}

Commercially available multigene assays such as Oncotype DX \citep{paik2004} and MammaPrint \citep{vantveer2002,cardoso2016} have demonstrated that transcriptomic information can meaningfully inform treatment decisions. These panels, however, are proprietary, fixed in composition, and return a single risk score without exposing the internal logic. At the other extreme, research models built on genome-wide features may achieve high discrimination but function as opaque ``black boxes'' whose predictions cannot be traced back to recognizable biological processes \citep{rudin2019}.

A growing body of work has applied explainable artificial intelligence (XAI) methods to cancer genomics. Chakraborty et al. \citep{chakraborty2021} used XGBoost with SHAP to relate immune cell composition to five-year breast cancer survival. Nkurunziza et al. \citep{nature2025shap} applied SHAP to genome-wide features for breast cancer diagnosis. Chen et al. \citep{chen2021pathway} constructed a sub-pathway activity matrix from KEGG annotations for survival stratification. Standard pathway-scoring methods such as single-sample gene set enrichment analysis (ssGSEA) \citep{barbie2009} and Gene Set Variation Analysis (GSVA) \citep{hanzelmann2013} are widely used but produce scores that are not directly optimized for supervised prediction and lack formal comparison with simpler aggregation strategies in prognostic contexts.

Despite this progress, no published study, to our knowledge, has combined (1)~manually curated hallmark pathway scores as the primary supervised feature set, (2)~both survival analysis and binary classification on the same features, (3)~formal benchmarking against ssGSEA, (4)~per-patient SHAP decomposition, (5)~controlled feature ablation, (6)~feature stability analysis, and (7)~external validation on a large independent cohort, all within a unified framework for breast cancer progression prediction.

\subsection{Contributions}

This study makes three principal contributions:

\begin{enumerate}
    \item \textbf{A biologically informed dimensionality-reduction strategy for prognosis.} We demonstrate that collapsing RNA-seq data into seven pathway scores, computed as simple mean z-scores, produces a compact feature representation that outperforms the standard ssGSEA enrichment baseline and carries independent prognostic signal both in survival analysis (concordance index) and binary classification (AUC-ROC).
    \item \textbf{A generalizable ML pipeline validated across cohorts and modeling paradigms.} The framework is evaluated through Cox proportional hazards, elastic net logistic regression, random forest, and gradient boosting, with systematic ablation and stability analyses, on two independent cohorts totaling 1,696 patients.
    \item \textbf{A patient-level interpretability framework using SHAP at pathway resolution.} Per-patient SHAP decompositions produce explanations grounded in recognized cancer hallmarks rather than opaque gene-level signals, with demonstrated stability across cross-validation folds (Spearman $\rho = 0.82$).
\end{enumerate}

\subsection{Scope and Non-Claims}

This work is a retrospective proof-of-concept study. The models described here are not intended as diagnostic tools and have not been validated in a prospective clinical setting. The framework is designed for research-grade risk stratification and as a foundation for future clinical decision support development. All analyses are observational; no causal claims regarding pathway activity and patient outcomes are made.


%======================================================================
\section{Methods}
\label{sec:methods}
%======================================================================

\subsection{Ethical Statement and Data Sources}

All data analyzed in this study are publicly available and fully de-identified. No new patient specimens were collected and no experimental procedures involving human subjects were performed. The study qualifies as ``Not Human Subjects Research'' under 45~CFR~\S46.102(e) and is exempt from Institutional Review Board review.

Two independent cohorts were used:

\textbf{Training cohort (TCGA-BRCA).} RNA-sequencing data (log$_2$ RSEM) and matched clinical annotations were obtained from The Cancer Genome Atlas Breast Invasive Carcinoma project \citep{tcga2012}. After excluding samples with incomplete survival data, 213 patients were retained.

\textbf{Validation cohort (GSE96058 / SCAN-B).} An independent set of 1,483 breast cancer samples with RNA-seq expression and clinical outcome data was drawn from the Sweden Cancerome Analysis Network, Breast initiative \citep{saal2015,brueffer2018}, accessed through the NCBI Gene Expression Omnibus (accession GSE96058). This cohort was reserved exclusively for validation; no GSE96058 data entered model training or hyperparameter tuning.

\subsection{Outcome Variables}

Two endpoints were defined. For \textbf{survival analysis}, the primary outcome was time to event (months from diagnosis to death or last follow-up), with event status indicating whether the patient died or experienced documented disease progression. For \textbf{binary classification}, a five-year endpoint was derived: patients who experienced an event within 60 months of diagnosis were labeled ``High Risk''; patients who survived beyond 60 months without documented progression were labeled ``Low Risk.'' Patients with follow-up shorter than five years and no recorded event were excluded.

\subsection{Expression Data Preprocessing}

Raw RNA-seq counts were log$_2$-transformed to stabilize variance. Within each cohort, expression values were z-score normalized across samples to reduce platform-specific and batch-related distributional shifts. Gene identifiers were mapped to HGNC-approved symbols.

\subsection{Biologically Informed Feature Engineering}

The central design choice of this study was to replace the raw expression matrix (more than 20,000 genes per sample) with seven pathway activity scores, each computed as the mean z-score normalized expression across a curated gene set. Gene lists were assembled from the MSigDB Hallmark collection \citep{liberzon2015} and refined using supporting cancer biology literature \citep{hanahan2011}. The seven pathways and their constituent genes are shown in Table~\ref{tab:pathways}.

\begin{table}[H]
\centering
\caption{Biological Pathway Definitions and Constituent Marker Genes}
\label{tab:pathways}
\small
\begin{tabular}{p{3cm} p{7.5cm} p{3cm}}
\toprule
\textbf{Pathway} & \textbf{Genes} & \textbf{Biological Role} \\
\midrule
Proliferation & \textit{MKI67, AURKA, BIRC5, CCNB1, MYBL2, PLK1, BUB1, CDC20, MCM2, PCNA} & Cell cycle activity \\
Estrogen Response & \textit{ESR1, PGR, FOXA1, GATA3, XBP1, CA12} & Hormone receptor signaling \\
Immune Response & \textit{CD8A, CD4, PDCD1, CD274, CTLA4, FOXP3, IL2RA} & Immune infiltration \\
Invasion / EMT & \textit{VIM, CDH2, FN1, MMP9, TWIST1, SNAI1, ZEB1} & Mesenchymal transition \\
Apoptosis & \textit{CASP3, CASP8, BAX, BAK1, BAD, TP53} & Programmed cell death \\
HER2 Signaling & \textit{ERBB2, GRB7, PGAP3, MIEN1, STARD3} & Receptor tyrosine kinase \\
Angiogenesis & \textit{VEGFA, KDR, FLT1, FGF2} & Neovascularization \\
\bottomrule
\end{tabular}
\end{table}

This aggregation serves three purposes. First, it reduces dimensionality by roughly three orders of magnitude, lowering overfitting risk. Second, it encodes prior biological knowledge directly into the feature space. Third, it aligns model features with concepts that oncologists already use when reasoning about tumor behavior. One engineered ratio feature, the proliferation-to-apoptosis ratio, was included to capture the net balance between cell division and programmed cell death.

For the GSE96058 validation cohort, nine standard clinical variables were additionally available: patient age at diagnosis, primary tumor size, lymph node status, estrogen receptor (ER) status, progesterone receptor (PgR) status, HER2 status, Ki-67 immunohistochemistry status, Nottingham Histological Grade (NHG), and PAM50 molecular subtype.

\subsection{Survival Analysis}

Cox proportional hazards models \citep{cox1972} with L2 regularization (penalizer = 0.1) were fitted on three feature configurations: pathway-only (8 features), clinical-only (9 features), and combined (17 features). The concordance index (C-index) was evaluated under stratified five-fold cross-validation with fixed random seeds. To visualize risk stratification, patients were divided at the median Cox-predicted partial hazard into high- and low-risk groups, and Kaplan--Meier survival curves were compared using the log-rank test.

\subsection{Binary Classification Models}

Three classifiers were compared using stratified five-fold cross-validation to preserve class proportions:

\begin{enumerate}
    \item \textbf{Elastic net logistic regression} (L1 + L2 penalty, saga solver, $l_1$-ratio = 0.5, $C = 1.0$), serving as a regularized linear baseline.
    \item \textbf{Random forest} (500 trees, maximum depth = 8, balanced class weights), an ensemble of decorrelated decision trees.
    \item \textbf{Gradient boosting} (500 estimators, learning rate = 0.03, maximum depth = 3), a sequential ensemble that iteratively corrects residual errors.
\end{enumerate}

The primary metric was AUC-ROC. Features were standardized within each fold (fit on training, transform on test). Class imbalance was handled through stratified fold construction and class-weight adjustments.

\subsection{Feature Ablation Study}

To isolate the contribution of pathway scores from clinical variables, three feature configurations were evaluated under identical five-fold cross-validation on GSE96058: pathway-only (8 features), clinical-only (9 features), and combined (17 features). Both Cox survival and all three binary classifiers were tested on each configuration.

\subsection{Pathway Scoring Baseline: ssGSEA}

To benchmark the mean z-score aggregation against an established pathway-scoring method, we implemented a rank-based single-sample gene set enrichment analysis (ssGSEA) baseline \citep{barbie2009}. For each sample, genes were ranked by expression value. An enrichment score was computed for each pathway gene set using the normalized rank-sum deviation from the expected null. The resulting ssGSEA scores were evaluated using the same classifiers and Cox model under identical cross-validation conditions, providing a controlled comparison of pathway aggregation strategies.

\subsection{Explainability Analysis}

SHAP (SHapley Additive exPlanations) values \citep{lundberg2017} were computed using the TreeExplainer algorithm for gradient boosting models on both the combined (17-feature) and pathway-only (8-feature) configurations. SHAP decomposes each prediction into additive per-feature contributions, providing both global importance rankings and patient-level explanations.

\subsection{Feature Importance Stability}

To assess whether SHAP-based feature rankings are stable or artifacts of particular data splits, SHAP importance was computed independently within each of the five cross-validation folds. Pairwise Spearman rank correlations between fold-specific feature importance vectors were computed, and the mean correlation was reported as a stability metric. This directly addresses the common concern that explainability rankings in omics settings are unreliable.

\subsection{Model Calibration}

Calibration curves were generated from out-of-fold predicted probabilities for all three binary classifiers on GSE96058 using the combined feature set. Calibration assesses whether a model's predicted probability of 0.7, for example, corresponds to an observed event rate of approximately 70\%.

\subsection{Code and Data Availability}

All preprocessing scripts, feature engineering pipelines, model training code, ablation analysis, and figure generation code are publicly available at: \url{https://github.com/suhaanthayyil/breast-cancer-prediction}. Raw data can be accessed through the GDC Data Portal (TCGA) and NCBI GEO (accession GSE96058).


%======================================================================
\section{Results}
\label{sec:results}
%======================================================================

\subsection{Cohort Characteristics}

The TCGA training cohort comprised 213 breast cancer patients with complete RNA-seq expression and survival annotation. The GSE96058 validation cohort included 1,483 patients from the SCAN-B initiative with clinical variables including age, tumor size, lymph node status, receptor status, Ki-67, histological grade, and PAM50 subtype.

\subsection{Survival Analysis}

Table~\ref{tab:cox} reports concordance indices for Cox proportional hazards models.

\begin{table}[H]
\centering
\caption{Cox Proportional Hazards Concordance Index}
\label{tab:cox}
\begin{tabular}{l c c}
\toprule
\textbf{Feature Set} & \textbf{C-index (5-fold CV)} & \textbf{C-index (full)} \\
\midrule
\multicolumn{3}{l}{\textit{GSE96058 ($n = 1{,}483$)}} \\
\quad Pathway Only (8 features) & 0.626 $\pm$ 0.030 & 0.640 \\
\quad Clinical Only (9 features) & 0.822 $\pm$ 0.030 & 0.824 \\
\quad Combined (17 features) & \textbf{0.827 $\pm$ 0.031} & \textbf{0.832} \\
\midrule
\multicolumn{3}{l}{\textit{TCGA ($n = 213$)}} \\
\quad Pathway Only (8 features) & --- & 0.590 \\
\bottomrule
\end{tabular}
\end{table}

The combined model achieved the highest C-index (0.827 cross-validated), with a pathway contribution of +0.005 over clinical features alone. The TCGA pathway-only C-index (0.590) was modest, consistent with the small sample size.

Kaplan--Meier analysis (Fig.~\ref{fig:km}) confirmed highly significant separation between Cox-predicted risk groups (log-rank $p < 5.7 \times 10^{-53}$). The high-risk group showed approximately 40\% lower survival probability at 80 months compared to the low-risk group.

\begin{figure}[H]
    \centering
    \includegraphics[width=0.75\textwidth]{fig_kaplan_meier.png}
    \caption{Kaplan--Meier survival curves for patients stratified by median Cox-predicted partial hazard (combined features, GSE96058). Shaded regions show 95\% confidence intervals. Log-rank $p < 5.7 \times 10^{-53}$.}
    \label{fig:km}
\end{figure}

\subsection{Binary Classification: TCGA (Pathway-Only)}

On TCGA (pathway-only features), the elastic net achieved AUC = 0.79 with accuracy = 0.77, outperforming both random forest (AUC = 0.67) and gradient boosting (AUC = 0.55). The linear model's advantage on this small, eight-feature dataset reflects the well-known bias--variance tradeoff: regularized linear methods outperform complex nonlinear models when the number of observations is limited (Fig.~\ref{fig:comparison}).

\begin{figure}[H]
    \centering
    \includegraphics[width=0.75\textwidth]{fig_comparison.png}
    \caption{Model performance comparison across datasets. AUC-ROC for each classifier on TCGA (blue; pathway-only) and GSE96058 (red; pathway + clinical). Dashed line at AUC = 0.75.}
    \label{fig:comparison}
\end{figure}

\subsection{Binary Classification: GSE96058 and Feature Ablation}

Table~\ref{tab:ablation} reports the full ablation results on GSE96058.

\begin{table}[H]
\centering
\caption{Feature Ablation: AUC-ROC on GSE96058 (Stratified 5-Fold CV)}
\label{tab:ablation}
\begin{tabular}{l c c c}
\toprule
\textbf{Feature Set} & \textbf{Elastic Net} & \textbf{Random Forest} & \textbf{Gradient Boosting} \\
\midrule
Pathway Only (8) & 0.641 $\pm$ 0.035 & 0.645 $\pm$ 0.030 & 0.633 $\pm$ 0.015 \\
Clinical Only (9) & 0.847 $\pm$ 0.033 & 0.849 $\pm$ 0.029 & 0.851 $\pm$ 0.027 \\
Combined (17) & \textbf{0.855 $\pm$ 0.031} & \textbf{0.856 $\pm$ 0.032} & 0.827 $\pm$ 0.025 \\
\midrule
$\Delta$ (Combined $-$ Clinical) & +0.008 & +0.007 & $-$0.024 \\
\bottomrule
\end{tabular}
\end{table}

Pathway-only models achieved AUC 0.633--0.645, demonstrating that curated pathway scores carry genuine prognostic signal from gene expression alone. Clinical-only models reached 0.847--0.851. The combined set yielded the highest AUC for elastic net (0.855) and random forest (0.856), with a consistent additive benefit. The gradient boosting decrease ($-$0.024) with 17 features likely reflects overfitting in the sequential boosting architecture.

\begin{figure}[H]
    \centering
    \includegraphics[width=\textwidth]{fig_ablation.png}
    \caption{Feature set ablation study on GSE96058. \textbf{(A)} AUC-ROC comparison across three feature configurations and three classifiers. Error bars: CV standard deviation. \textbf{(B)} ROC curves for the elastic net classifier under each configuration.}
    \label{fig:ablation}
\end{figure}

\subsection{ssGSEA Baseline Comparison}

Table~\ref{tab:ssgsea} compares the mean z-score aggregation against the ssGSEA baseline using pathway-only features.

\begin{table}[H]
\centering
\caption{Mean Z-score vs. ssGSEA Pathway Scoring (Pathway-Only Features)}
\label{tab:ssgsea}
\begin{tabular}{l c c c}
\toprule
\textbf{Model} & \textbf{Mean Z-score} & \textbf{ssGSEA} & \textbf{$\Delta$} \\
\midrule
Elastic Net (AUC) & 0.641 & 0.631 & +0.010 \\
Random Forest (AUC) & 0.645 & 0.622 & +0.023 \\
Gradient Boosting (AUC) & 0.633 & 0.595 & +0.038 \\
Cox PH (C-index) & 0.640 & 0.636 & +0.004 \\
\bottomrule
\end{tabular}
\end{table}

The mean z-score aggregation consistently outperformed ssGSEA across all four modeling approaches, with improvements ranging from +0.004 (Cox C-index) to +0.038 (gradient boosting AUC). This finding has practical implications discussed in Section~\ref{sec:discussion}.

\subsection{ROC Curves}

\begin{figure}[H]
    \centering
    \includegraphics[width=0.85\textwidth]{fig_roc.png}
    \caption{ROC curves. \textbf{Left:} TCGA, elastic net (AUC = 0.79). \textbf{Right:} GSE96058 five-fold CV, gradient boosting (AUC = 0.83).}
    \label{fig:roc}
\end{figure}

\subsection{SHAP Feature Importance and Beeswarm Analysis}

The SHAP beeswarm plot (Fig.~\ref{fig:shap_combined}) reveals the feature attributions for the combined gradient boosting model on GSE96058. Ki-67 immunohistochemistry status and patient age at diagnosis were the two most influential features by a substantial margin, followed by the estrogen response, immune response, and proliferation pathway scores. Binary receptor status variables (ER, PgR, HER2) contributed relatively little beyond what the pathway scores already captured.

Several biological patterns emerge from the directional effects.

\textit{Ki-67 and proliferation as risk amplifiers.} High Ki-67 status was consistently pushed toward positive SHAP values (higher risk). The proliferation pathway score showed a similar, though weaker, pattern. This is expected: Ki-67 is a well-established marker of mitotic activity \citep{dowsett2011}.

\textit{Estrogen response as a protective signal.} High estrogen response pathway scores were associated with negative SHAP values, matching the clinical reality that ER-positive tumors generally carry a more favorable prognosis.

\textit{Immune response and invasion/EMT.} Elevated immune response showed a mild protective association, consistent with evidence that cytotoxic T-cell infiltration improves outcomes \citep{salgado2015}. Higher invasion/EMT activity pushed predictions toward elevated risk.

\textit{Age at diagnosis.} Older age was associated with higher predicted risk, with wide SHAP spread suggesting context-dependent contributions.

\begin{figure}[H]
    \centering
    \includegraphics[width=0.85\textwidth]{fig_shap_beeswarm.png}
    \caption{SHAP beeswarm for the combined gradient boosting model on GSE96058 ($n = 1{,}483$). Each dot represents one patient. Features ranked by mean $|$SHAP$|$. Color: red = high feature value, blue = low. Ki-67 and age dominate; estrogen response ranks third.}
    \label{fig:shap_combined}
\end{figure}

In the pathway-only model (Fig.~\ref{fig:shap_pw}), estrogen response emerged as the dominant contributor, followed by the proliferation-to-apoptosis ratio and immune response. Without clinical features to absorb signal, the model fully exposes the molecular drivers of risk.

\begin{figure}[H]
    \centering
    \includegraphics[width=0.85\textwidth]{fig_shap_pathway_only.png}
    \caption{SHAP beeswarm for the pathway-only gradient boosting model on GSE96058 (8 features). Estrogen response dominates when clinical proxies are removed.}
    \label{fig:shap_pw}
\end{figure}

\subsection{Patient-Level Explanations}

SHAP waterfall plots for representative patients (Fig.~\ref{fig:waterfalls}) demonstrate that the model's reasoning varies across patients in recognizable ways. For the high-risk patient, Ki-67 (+4.97 SHAP units) and age (+1.48) dominate. For the low-risk patient, younger age ($-$0.92), low proliferation ($-$0.46), high apoptosis ($-$0.43), and low Ki-67 ($-$0.43) drive the prediction.

\begin{figure}[H]
    \centering
    \begin{subfigure}[t]{0.85\textwidth}
        \centering
        \includegraphics[width=\textwidth]{fig_waterfall_high.png}
        \caption{High-risk patient: $f(x) = 6.339$}
    \end{subfigure}
    \vspace{1em}
    \begin{subfigure}[t]{0.85\textwidth}
        \centering
        \includegraphics[width=\textwidth]{fig_waterfall_low.png}
        \caption{Low-risk patient: $f(x) = -5.711$}
    \end{subfigure}
    \caption{SHAP waterfall plots for representative patients from the GSE96058 combined model.}
    \label{fig:waterfalls}
\end{figure}

\subsection{Feature Importance Stability}

SHAP feature importance rankings were computed independently within each of the five cross-validation folds. The mean pairwise Spearman rank correlation was $\rho = 0.82$ (Fig.~\ref{fig:stability}), indicating high stability. Ki-67 and age ranked in the top 2 in all five folds. Pathway features showed somewhat greater rank variability but remained consistently in the top half of the importance hierarchy.

\begin{figure}[H]
    \centering
    \includegraphics[width=\textwidth]{fig_stability.png}
    \caption{Feature importance stability. \textbf{(A)} Heatmap of mean $|$SHAP$|$ per fold (darker = more important). \textbf{(B)} Rank stability as boxplots; narrow boxes indicate consistent ranking. Mean Spearman $\rho = 0.82$.}
    \label{fig:stability}
\end{figure}

\subsection{Model Calibration}

Calibration curves (Fig.~\ref{fig:cal}) reveal that random forest produced the best-calibrated predictions, closely tracking the diagonal. The elastic net exhibited systematic underestimation of risk in the mid-range, while gradient boosting showed step-like behavior. For clinical deployment, these results suggest random forest offers the best combination of discrimination and calibration.

\begin{figure}[H]
    \centering
    \includegraphics[width=0.7\textwidth]{fig_calibration.png}
    \caption{Calibration curves for three classifiers on GSE96058 (combined features, out-of-fold predictions). Dashed diagonal: perfect calibration. Random forest tracks most closely.}
    \label{fig:cal}
\end{figure}


%======================================================================
\section{Discussion}
\label{sec:discussion}
%======================================================================

\subsection{Summary of Findings}

This study presents a biologically informed explainable machine learning framework for breast cancer progression risk stratification. The framework was evaluated through six complementary analyses: (1) Cox survival modeling with concordance index and Kaplan--Meier visualization, (2) binary classification with three classifier families, (3) controlled feature ablation, (4) ssGSEA baseline comparison, (5) SHAP explainability with formal stability analysis, and (6) model calibration, all validated on an external cohort of 1,483 patients.

\subsection{Survival Analysis Validates the Binary Endpoint}

The inclusion of Cox proportional hazards modeling addresses a fundamental limitation of binary-only analysis: it confirms that the features predict time-to-event outcomes, not merely a dichotomized label. The combined C-index (0.827) is competitive with published Cox models in breast cancer that use larger feature sets \citep{chen2021pathway}. The Kaplan--Meier curves demonstrate visually striking risk-group separation (log-rank $p < 10^{-52}$), providing evidence that is intuitive to clinical audiences. Importantly, the survival analysis results closely mirror the binary classification results: in both paradigms, pathway features carry independent signal that provides a consistent additive benefit when combined with clinical variables.

\subsection{Mean Z-score Outperforms ssGSEA for Curated Gene Sets}

A key finding is that the simple mean z-score pathway aggregation consistently outperformed the ssGSEA baseline across all four modeling approaches. This is not paradoxical. ssGSEA was designed for genome-scale gene sets (hundreds of genes) where rank-based enrichment captures subtle distributional shifts. For small, curated gene sets of 4--10 genes, a simple average of standardized values preserves more discriminative information than rank-based statistics, which lose magnitude information during the ranking step.

This finding has practical implications: the simpler method is not only more performant in this context but also more transparent (a mean z-score is immediately interpretable) and more computationally efficient. It should not be generalized to large gene sets, where ssGSEA and GSVA remain appropriate.

\subsection{Interpreting the Ablation Results}

Pathway scores alone achieve moderate discrimination (AUC $\approx$ 0.64, C-index $\approx$ 0.63). When combined with clinical features, the improvement is modest in magnitude (+0.007 AUC, +0.005 C-index for elastic net/random forest) but consistent. The modest magnitude reflects partial redundancy between pathway scores and clinical variables: Ki-67 is a single-gene proxy for our ten-gene proliferation pathway; ER status binarizes our continuous estrogen response score.

The pathway scores' primary contribution is therefore not raw discrimination but \textit{representational richness}. Where clinical variables report binary categories, pathway scores provide continuous biological axes. This enables the per-patient SHAP decompositions that constitute the central interpretability contribution of this work.

\subsection{Feature Stability Supports Generalizability}

The high cross-fold rank stability (Spearman $\rho = 0.82$) indicates that the model's learned importance hierarchy is robust to data perturbation, not an artifact of particular data splits. Ki-67 and age consistently ranked in the top 2. Pathway features showed moderate rank variability consistent with their smaller effect sizes. This stability analysis directly addresses the common concern that SHAP rankings in omics settings are unreliable.

\subsection{Calibration and Clinical Utility}

The calibration analysis reveals that random forest offers the best trade-off between discrimination (AUC = 0.856) and calibration fidelity. For any clinical deployment, calibration is arguably as important as discrimination: a poorly calibrated model that assigns a ``70\% risk'' when the true risk is 40\% could lead to inappropriate treatment escalation.

\subsection{Relationship to Existing Work}

Chakraborty et al. \citep{chakraborty2021} demonstrated that SHAP could recover biologically plausible attributions using immune cell fractions. Our approach captures seven distinct cancer hallmarks rather than immune composition alone. The finding that our immune response pathway score behaves similarly to their immune cell fractions provides convergent validation.

Chen et al. \citep{chen2021pathway} used KEGG sub-pathway features for survival stratification, but their pipeline relied on unsupervised consensus clustering without per-patient explainability. Our framework extends this line of work by pairing pathway scores with supervised classifiers and post-hoc explanation, adding ablation, stability analysis, and ssGSEA benchmarking.

\subsection{Limitations}

The TCGA training cohort ($n = 213$) is small by machine learning standards. Pathway scores computed as simple gene-set means do not account for intra-pathway regulatory structure; more sophisticated aggregation (ssGSEA, network-informed scoring) may capture additional signal with larger gene sets. The modest pathway contribution (+0.007 AUC) reflects redundancy with clinical proxies; datasets lacking clinical annotation would likely show a larger effect. All analyses are retrospective. Model calibration and clinical utility in prospective, multi-institutional settings remain to be demonstrated.

\subsection{Future Directions}

Several extensions follow naturally: (1) time-to-event modeling with random survival forests and DeepSurv; (2) application to other solid tumor types; (3) assessment of pathway scores as predictors of differential treatment response; (4) federated learning across institutions; and (5) cross-cohort transfer analysis (training on TCGA, testing on METABRIC) to further test generalizability.


%======================================================================
\section{Conclusion}
\label{sec:conclusion}
%======================================================================

We present a biologically informed explainable machine learning framework for breast cancer progression risk stratification from tumor gene expression. Through survival analysis (Cox C-index 0.827, Kaplan--Meier log-rank $p < 10^{-52}$), binary classification (AUC 0.856), feature ablation, ssGSEA baseline comparison, SHAP explainability with stability assessment (Spearman $\rho = 0.82$), and model calibration, we demonstrate that curated pathway scores carry independent prognostic signal, outperform standard enrichment methods for small gene sets, and provide biologically grounded per-patient explanations. The framework's transparency, stability, and generalizability position it as a foundation for interpretable clinical decision support in breast cancer genomics.


%======================================================================
\bibliographystyle{unsrtnat}
\begin{thebibliography}{99}

\bibitem[Barbie et al.(2009)]{barbie2009}
D.~A.~Barbie, P.~Tamayo, J.~S.~Boehm, et~al.
\newblock Systematic RNA interference reveals that oncogenic KRAS-driven cancers require TBK1.
\newblock \textit{Nature}, 462(7269):108--112, 2009.

\bibitem[Brueffer et al.(2018)]{brueffer2018}
C.~Brueffer, J.~Vallon-Christersson, D.~Grabau, et~al.
\newblock Clinical value of RNA sequencing-based classifiers for prediction of the five conventional breast cancer biomarkers.
\newblock \textit{JCO Precision Oncology}, 2:1--18, 2018.

\bibitem[Cardoso et al.(2016)]{cardoso2016}
F.~Cardoso, L.~J.~van~'t~Veer, J.~Bogaerts, et~al.
\newblock 70-gene signature as an aid to treatment decisions in early-stage breast cancer.
\newblock \textit{New England Journal of Medicine}, 375(8):717--729, 2016.

\bibitem[Chakraborty et al.(2021)]{chakraborty2021}
D.~Chakraborty, C.~Ivan, P.~Amero, et~al.
\newblock Explainable artificial intelligence reveals novel insight into tumor microenvironment conditions linked with better prognosis in patients with breast cancer.
\newblock \textit{Cancers}, 13(14):3450, 2021.

\bibitem[Chen et al.(2021)]{chen2021pathway}
L.~Chen, Y.~Wang, J.~Zhao, et~al.
\newblock Identification of pathway-based biomarkers with crosstalk analysis for overall survival risk prediction in breast cancer.
\newblock \textit{Frontiers in Genetics}, 12:689715, 2021.

\bibitem[Cox(1972)]{cox1972}
D.~R.~Cox.
\newblock Regression models and life-tables.
\newblock \textit{Journal of the Royal Statistical Society: Series B}, 34(2):187--220, 1972.

\bibitem[Dowsett et al.(2011)]{dowsett2011}
M.~Dowsett, T.~O.~Nielsen, R.~A'Hern, et~al.
\newblock Assessment of Ki67 in breast cancer: Recommendations from the International Ki67 in Breast Cancer Working Group.
\newblock \textit{Journal of the National Cancer Institute}, 103(22):1656--1664, 2011.

\bibitem[Hanahan and Weinberg(2011)]{hanahan2011}
D.~Hanahan and R.~A.~Weinberg.
\newblock Hallmarks of cancer: The next generation.
\newblock \textit{Cell}, 144(5):646--674, 2011.

\bibitem[H\"{a}nzelmann et al.(2013)]{hanzelmann2013}
S.~H\"{a}nzelmann, R.~Castelo, and J.~Guinney.
\newblock GSVA: Gene set variation analysis for microarray and RNA-seq data.
\newblock \textit{BMC Bioinformatics}, 14(1):7, 2013.

\bibitem[Liberzon et al.(2015)]{liberzon2015}
A.~Liberzon, C.~Birger, H.~Thorvaldsd\'{o}ttir, et~al.
\newblock The Molecular Signatures Database (MSigDB) Hallmark gene set collection.
\newblock \textit{Cell Systems}, 1(6):417--425, 2015.

\bibitem[Lundberg and Lee(2017)]{lundberg2017}
S.~M.~Lundberg and S.-I.~Lee.
\newblock A unified approach to interpreting model predictions.
\newblock In \textit{Advances in Neural Information Processing Systems}, 30:4765--4774, 2017.

\bibitem[Nkurunziza et al.(2025)]{nature2025shap}
H.~Nkurunziza, M.~Schmid, and G.~Casalicchio.
\newblock Breast cancer prediction based on gene expression data using interpretable machine learning techniques.
\newblock \textit{Scientific Reports}, 15:7747, 2025.

\bibitem[Paik et al.(2004)]{paik2004}
S.~Paik, S.~Shak, G.~Tang, et~al.
\newblock A multigene assay to predict recurrence of tamoxifen-treated, node-negative breast cancer.
\newblock \textit{New England Journal of Medicine}, 351(27):2817--2826, 2004.

\bibitem[Rudin(2019)]{rudin2019}
C.~Rudin.
\newblock Stop explaining black box machine learning models for high stakes decisions and use interpretable models instead.
\newblock \textit{Nature Machine Intelligence}, 1(5):206--215, 2019.

\bibitem[Saal et al.(2015)]{saal2015}
L.~H.~Saal, J.~Vallon-Christersson, J.~Hakkinen, et~al.
\newblock The Sweden Cancerome Analysis Network--Breast (SCAN-B) initiative.
\newblock \textit{Genome Medicine}, 7(1):20, 2015.

\bibitem[Salgado et al.(2015)]{salgado2015}
R.~Salgado, C.~Denkert, S.~Demaria, et~al.
\newblock The evaluation of tumor-infiltrating lymphocytes (TILs) in breast cancer.
\newblock \textit{Annals of Oncology}, 26(2):259--271, 2015.

\bibitem[van~'t~Veer et al.(2002)]{vantveer2002}
L.~J.~van~'t~Veer, H.~Dai, M.~J.~van~de~Vijver, et~al.
\newblock Gene expression profiling predicts clinical outcome of breast cancer.
\newblock \textit{Nature}, 415(6871):530--536, 2002.

\bibitem[TCGA(2012)]{tcga2012}
The Cancer Genome Atlas Network.
\newblock Comprehensive molecular portraits of human breast tumours.
\newblock \textit{Nature}, 490(7418):61--70, 2012.

\end{thebibliography}

\end{document}
